\topic{rigid-body}{Torque-free rigid-body dynamics on SO(3)}
\study{rigid-body}
Reading anchors: Wie, second edition, Section 6.4, printed p.~359 / PDF
p.~377, and Section 6.5, printed p.~363 / PDF p.~381 \cite{wie}.
These are exact indexed mappings. Section 6.2 concerns inertia, not Euler's
equations. Ernest's DGDT and quaternion notes provide the geometric bridge.

\subsection{Inertia and Euler's equations}
Let $R$ map body to inertial coordinates and let
$I=\operatorname{diag}(I_1,I_2,I_3)$ be a constant positive-definite inertia
tensor in body principal axes. Body angular momentum is $\Pi=I\omega$;
inertial angular momentum is $L=R\Pi$. In absence of external torque,
\[
 0=\dot L=R(\dot\Pi+\omega\times\Pi),\qquad
 I\dot\omega+\omega\times I\omega=0.
\]
Thus $\dot\omega_1=(I_2-I_3)\omega_2\omega_3/I_1$, with cyclic analogues.
Multiplying by $\omega$ gives
$\dot E=\omega\cdot I\dot\omega=-\omega\cdot(\omega\times\Pi)=0$.
Likewise $\tfrac{d}{dt}|\Pi|^2=-2\Pi\cdot(\omega\times\Pi)=0$.
These conservation statements require no choice of Euler-angle chart.

\subsection{A solvable axisymmetric case}
For $I_1=I_2=I_\perp$, the third component is constant, $\omega_3=c$.
Writing $\nu=(I_3-I_\perp)c/I_\perp$, the transverse equations are
$\dot\omega_1=-\nu\omega_2$, $\dot\omega_2=\nu\omega_1$, hence
\[
 \begin{pmatrix}\omega_1(t)\\\omega_2(t)\end{pmatrix}
 =\begin{pmatrix}\cos\nu t&-\sin\nu t\\\sin\nu t&\cos\nu t\end{pmatrix}
 \begin{pmatrix}a\\b\end{pmatrix}.
\]
The worked data are $(I_1,I_2,I_3)=(2,2,3)\,\mathrm{kg\,m^2}$,
$(a,b,c)=(0.4,0.2,0.7)\,\mathrm{rad/s}$, $R(0)=1$, and $t=10\,\mathrm{s}$.
These are explicitly chosen synthetic inputs, not textbook or spacecraft data.

\subsection{What the executable check establishes}
The teaching implementation integrates $(\omega,R)$ using classical RK4,
with $\dot R=R\widehat\omega$. It compares body spin to the analytic solution
at 200 and 400 steps, requiring an error ratio between 14 and 18.
At the finer resolution it requires maximum component error in $R^TR-1$
below $10^{-7}$, determinant error below $10^{-7}$, energy error below
$10^{-8}\,\mathrm{J}$, and inertial momentum error below
$10^{-7}\,\mathrm{kg\,m^2/s}$.

Ordinary RK4 does not preserve $\SO(3)$ exactly; the rotation constraints are
measured errors here. A future Lie-group integrator should preserve membership
by construction and still face independent accuracy checks. Sage proves the
energy and momentum-norm derivative identities symbolically for nonzero
principal moments. Neither those identities nor this one axisymmetric test
constitutes verification of arbitrary torque models.
